% ID: SEL-D0028-OB-E02
% ESTADO: cotejada. Segunda pasada visual completa y borrador compilado revisados.
% Fuente: Selectividad/matematicas_aplicadas/suplente junio/Examen-Matemáticas Aplicadas a las Ciencias Sociales.pdf
% SHA-256: a10146fce64a17e21450973c88d66d43571619781f76bc345b5a4b3dee67a16b
% Página: 2; opción B; ejercicio original 2.

Dada la función
\[
f(x)=\begin{cases}
\dfrac{2x+1}{1-2x} & \text{si }x<0\\[4pt]
x^2-x-a & \text{si }x\geq0
\end{cases}
\]
\begin{enumerate}[label=\alph*)]
\item \textbf{(1.3 puntos)} Obtenga el valor de $a$ para que la función sea continua en $x=0$. Para ese valor de $a$, \textquestiondown{}sería derivable en $x=0$?
\item \textbf{(1.2 puntos)} Para $a=2$, estudie su monotonía y extremos relativos.
\end{enumerate}
